\section{Introdução}

De acordo com \citet{bispo2020tecnologias}: 

\begin{quote}
Embora Informática na Educação (IE) e Educação de Computação (EC) partilhem muitos elementos metodológicos, epistemológicos, contextos de pesquisa e em alguns casos até objetivos de pesquisa convergentes, a EC distingue-se como área
de pesquisa específica por sua \emph{raison d’être}: a busca do entendimento profundo de fenômenos complexos e processos envolvidos no ensino e aprendizado de computação \cite{malmi2019computing}. Além disso, enquanto as TEC e a IE pressupõem o emprego de um artefato tecnológico no processo de ensino e/ou aprendizagem, a EC pode utilizar-se ou não de tais artefatos para o ensino dos processos de computação. Assim como a Ciência da Computação ainda pode ser considerada uma ciência relativamente nova e em desenvolvimento, a EC ainda encontra-se em um estado de maturação \cite{malmi2019computing}. Ainda que a maioria das teorias utilizadas para o ensino de computação seja adaptada de outras áreas de tecnologia (como a engenharia, matemática ou ciências de forma geral), ou de teorias de educação ou psicologia da educação de forma mais ampla \cite{malmi2014theoretical}, já é possível observar o advento de teorias e práticas pedagógicas próprias \cite{malmi2019computing} com um crescente suporte empírico \cite{malmi2014theoretical}
\end{quote}


\lipsum[3-10]